
\documentclass[11pt]{article}

\usepackage[margin=1in]{geometry}
\usepackage{amsmath,amssymb}
\usepackage{listings}
\usepackage{xcolor}
\usepackage{tikz}
\usepackage{hyperref}
\usetikzlibrary{shapes.geometric}

\definecolor{backcolour}{rgb}{0.95,0.95,0.92}

\lstset{
    backgroundcolor=\color{backcolour},
    basicstyle=\ttfamily\small,
    breaklines=true,
    frame=single
}

\title{CPSC 250 \\ Branching and Flow Control}
\author{Edward Brash}
\date{}

\begin{document}

\maketitle

\section{Introduction}

One of the most important ideas in programming is the ability for a program to make decisions.

Programs are not simply sequences of instructions executed line-by-line forever.

Instead, programs often need to:

\begin{itemize}
    \item make decisions,
    \item repeat actions,
    \item and alter execution flow dynamically.
\end{itemize}

This is called:

\[
\textbf{flow control}
\]

\section{Branching}

Branching occurs whenever a program must decide between multiple execution paths.

Conceptually:

\begin{center}
\begin{tikzpicture}[scale=1.2]

% Nodes
\node (start) at (0,2) {};
\node[draw, diamond, aspect=2] (decision) at (0,0) {Condition?};

\node (yes) at (-2,-2) {Yes Branch};
\node (no) at (2,-2) {No Branch};

% Arrows
\draw[->] (start) -- (decision);

\draw[->] (decision) -- node[left] {Y} (yes);
\draw[->] (decision) -- node[right] {N} (no);

\end{tikzpicture}
\end{center}


\section{The if-else Structure}

The most general branching structure in Python is:

\begin{lstlisting}[language=Python]
if condition:
    # do something
else:
    # do something else
\end{lstlisting}

Sometimes the \texttt{else} branch simply does nothing.

In that case:

\begin{lstlisting}[language=Python]
if condition:
    # do something
\end{lstlisting}

is really just shorthand for:

\begin{lstlisting}[language=Python]
if condition:
    # do something
else:
    pass
\end{lstlisting}

\section{Conditions Evaluate to Booleans}

A condition must evaluate to a Boolean value:

\[
\texttt{True}
\]

or

\[
\texttt{False}
\]

Example:

\begin{lstlisting}[language=Python]
a == 4
\end{lstlisting}

This expression evaluates to:

\begin{itemize}
    \item \texttt{True} if \texttt{a} equals 4
    \item \texttt{False} otherwise
\end{itemize}

\section{Comparison vs Assignment}

A very common mistake for beginners is confusing:

\begin{lstlisting}[language=Python]
==
\end{lstlisting}

with:

\begin{lstlisting}[language=Python]
=
\end{lstlisting}

These are completely different operations.

\subsection{Logical Comparison}

\begin{lstlisting}[language=Python]
a == 4
\end{lstlisting}

asks:

\begin{quote}
``Is the current value of \texttt{a} equal to 4?''
\end{quote}

\subsection{Assignment}

\begin{lstlisting}[language=Python]
a = 4
\end{lstlisting}

means:

\begin{quote}
``Store the value 4 into the memory location associated with \texttt{a}.''
\end{quote}

\section{Loops}

Loops allow a section of code to execute repeatedly.

There are two major loop structures in Python:

\begin{enumerate}
    \item \texttt{for} loops
    \item \texttt{while} loops
\end{enumerate}

\section{For Loops}

A \texttt{for} loop is useful when you know approximately how many times you want to execute a loop.

Example:

\begin{lstlisting}[language=Python]
for i in range(10):
    print(i)
\end{lstlisting}

This loop executes 10 times:

\[
i = 0,1,2,\ldots,9
\]

\subsection{Looping Through Collections}

\begin{lstlisting}[language=Python]
l = ['a', 'b', 'c']

for item in l:
    print(item)
\end{lstlisting}

Output:

\begin{lstlisting}
a
b
c
\end{lstlisting}

\section{While Loops}

A \texttt{while} loop is useful when you do \emph{not} know in advance how many iterations are required.

General structure:

\begin{lstlisting}[language=Python]
while condition:
    # do stuff
\end{lstlisting}

The loop continues executing while the condition remains true.

Example:

\begin{lstlisting}[language=Python]
a = 0

while a < 4:
    print(a)
    a += 1
\end{lstlisting}

The loop terminates once the condition becomes false.

\section{Infinite Loops}

One danger with \texttt{while} loops is accidentally creating infinite loops.

Example:

\begin{lstlisting}[language=Python]
while True:
    print("Oops")
\end{lstlisting}

This loop never terminates unless interrupted externally.

\section{Break Statements}

Sometimes loops contain multiple conditions for termination.

Example problem:

\begin{quote}
Loop through a list until reaching the value \texttt{"end"}.
\end{quote}

\begin{lstlisting}[language=Python]
l = ["bob", "alice", "fred", "end", "jane"]

for item in l:

    if item == "end":
        break

    print(item)
\end{lstlisting}

Output:

\begin{lstlisting}
bob
alice
fred
\end{lstlisting}

The \verb|break| statement immediately exits the loop.

\section{Key Takeaways}

\begin{itemize}
    \item Branching controls program decision-making.

    \item Conditions evaluate to Boolean values.

    \item \verb|==| performs comparison.

    \item \verb|=| performs assignment.

    \item \verb|for| loops are ideal for known iteration counts.

    \item \verb|while| loops are ideal for condition-based iteration.

    \item \verb|break| immediately exits a loop.
\end{itemize}

\end{document}
